\section{Theoretical Interpretation}\label{sec:theoretical}

We provide an interpretation of intrinsic token weighting as a variance-control and
credit-redistribution mechanism for policy gradient estimation in language models. Our
goal is not to derive new convergence guarantees, but to clarify how token weighting
relates to advantage estimation and why it can be effective without learning value
functions.

\subsection{Token Weighting as Credit Redistribution}

Consider the standard REINFORCE estimator with a batch-level baseline:
\begin{equation}
g = (R - b)\sum_{t=1}^{T} \nabla_\theta \log \pi_\theta(y_t \mid y_{<t}, x).
\end{equation}

This estimator assigns equal credit to all tokens in a trajectory. Introducing token
weights yields:
\begin{equation}
g_w = (R - b)\sum_{t=1}^{T} w_t(y)\, \nabla_\theta \log \pi_\theta(y_t \mid y_{<t}, x),
\end{equation}
where $\sum_t w_t = 1$.

Importantly, this transformation does not change the reward signal itself; it
redistributes how the trajectory-level signal is attributed to individual decisions.
When $w_t$ depends only on quantities available at sampling time (e.g.,
log-probabilities or entropy), the estimator remains on-policy and does not require
learning additional functions.

\subsection{Relationship to Advantage Estimation}

Advantage-based methods can be interpreted as implicitly defining token-level weights
via a learned value function:
\begin{equation}
A_t = R - V(y_{<t}, x).
\end{equation}
In this view, the contribution of each token is scaled according to how much the
observed outcome deviates from an estimated expectation conditioned on the prefix.

However, this interpretation relies on the existence of a meaningful value function over
prefixes. In language generation, prefixes are non-Markov and do not form a stable state
space, making $V(y_{<t}, x)$ difficult to estimate and prone to bias. Token weighting
offers an alternative: rather than estimating expected returns from prefixes, it
emphasizes tokens based on intrinsic measures of decision salience, such as uncertainty
or information gain.

\subsection{Connection to Control Variates}

From a statistical perspective, subtracting a baseline $b$ is a form of control variate
that reduces variance without affecting unbiasedness. Token weighting can be viewed as a
complementary mechanism that reshapes the contribution of individual score-function
terms:
\begin{equation}
\nabla_\theta \log \pi_\theta(y) = \sum_{t=1}^{T} \nabla_\theta \log \pi_\theta(y_t \mid y_{<t}, x).
\end{equation}

Uniform weighting treats all score-function terms equally, regardless of their variance
or relevance. Intrinsic weighting schemes effectively reweight these terms,
down-weighting low-variance or low-impact contributions (e.g., high-probability
continuation tokens) and emphasizing high-variance or high-impact decisions (e.g.,
low-probability or uncertainty-reducing tokens). While this reweighting introduces bias
relative to the uniform estimator, it can substantially reduce variance, leading to
improved optimization dynamics in practice.

\subsection{Interpretation of Specific Weighting Schemes}

\paragraph{Surprisal Weighting.}
Surprisal-based weights emphasize tokens with low conditional probability under the
current policy. These tokens correspond to decisions where small parameter changes can
induce large changes in likelihood, and therefore often dominate gradient variance.
Weighting by surprisal concentrates learning signal on such high-sensitivity decisions.

\paragraph{Entropy-Reduction Weighting.}
Entropy-reduction weighting emphasizes tokens that sharply reduce predictive entropy.
These tokens correspond to commitment points in generation, where the model transitions
from a diffuse distribution over continuations to a more concentrated one. This mirrors
the role of branching points in tree-based reasoning methods, but is computed locally
from the model's predictive distribution without explicit tree construction.

\subsection{Bias--Variance Tradeoff}

Both advantage estimation and intrinsic token weighting introduce bias in exchange for
variance reduction. Advantage estimation does so by relying on a learned value function,
whose bias can be substantial when the underlying state abstraction is weak. Intrinsic
token weighting introduces bias by reshaping the contribution of token-level gradients
based on heuristic salience measures. The empirical question is therefore not whether
bias is introduced, but whether the bias--variance tradeoff is favorable.

This motivates the empirical comparison implemented in the repository: if intrinsic
token weighting improves accuracy, pass@$k$, or optimization behavior under matched
settings, then the bias it introduces may be more benign than the bias induced by
ill-defined value targets over text prefixes. The point is not that weighting is
unbiased, but that its inductive bias may be better aligned with token-level credit
assignment in language generation.

\subsection{Signal Degeneration Under Low-Rank Adaptation}\label{sec:theory-collapse}

The preceding analysis assumes that the intrinsic weights
$\{w_t\}_{t=1}^T$ are non-degenerate: if every scheme becomes uniform or
spuriously sparse for reasons unrelated to the learned adapter, then token
weighting cannot provide meaningful credit redistribution over a uniform
baseline. We now formalize when this degeneracy occurs under LoRA.

\paragraph{Setup.}
Let $\pi_{\mathrm{ref}}$ be a frozen base model with last-layer hidden states
$h_t^{\mathrm{base}} \in \mathbb{R}^d$ and a LoRA adapter whose net
contribution at position $t$ is $\Delta h_t \in \mathbb{R}^d$, giving
adapted hidden states $h_t^{\mathrm{adapted}} = h_t^{\mathrm{base}} + \Delta
h_t$ and logits $z_t^{\theta} = W_{\mathrm{lm}} h_t^{\mathrm{adapted}}$.
Let $\beta = \max_t \|\Delta h_t\|_2$ denote the maximum per-position
adapter contribution and $L = \|W_{\mathrm{lm}}\|_{\mathrm{op}}$ its
logit-space Lipschitz constant. Write $G(\cdot)$ for the Gini coefficient
of a non-negative vector and $\mathrm{EffN}(w)/T = 1 / (T \sum_t w_t^2)$ for
the normalized effective-token count. When a normalization uses an
$\varepsilon$ floor, the statements below are interpreted in the regime where
the raw salience scores dominate $\varepsilon$; if fixed $\varepsilon$
dominates all scores, every scheme reverts to uniform by construction.

\begin{proposition}[Degeneration of surprisal and entropy weighting]
\label{prop:surprisal-collapse}
Let $\alpha_t^{\mathrm{surp}} = -\log \pi_\theta(y_t \mid y_{<t}, x)$ and
$\alpha_t^{\mathrm{surp,ref}} = -\log \pi_{\mathrm{ref}}(y_t \mid y_{<t}, x)$,
with corresponding normalized weights $w^{\mathrm{surp}}$ and
$w^{\mathrm{surp,ref}}$. Then for every trajectory,
\[
\lVert w^{\mathrm{surp}} - w^{\mathrm{surp,ref}} \rVert_1
\le \frac{4 L \beta}{\sum_t \alpha_t^{\mathrm{surp,ref}}}.
\]
In particular, the adapted surprisal weights remain close to the reference
model's surprisal profile rather than becoming an adapter-specific credit
signal. If that reference profile is close to uniform, the weights approach
$1/T$; if it is highly uneven, the weights can instead remain spuriously
concentrated for reasons inherited from the base model. The analogous
statement holds for the entropy-reduction score.
\end{proposition}

\noindent\emph{Proof sketch.}
The softmax Jacobian is $2$-Lipschitz in the log-space norm, so the per-token
log-probabilities differ by at most $2 L \beta$ between the adapted and base
model. The projection onto the simplex via within-trajectory normalization is
$1$-Lipschitz in $\ell_1$ up to a factor of $2 / \sum_t \alpha_t^{\mathrm{ref}}$;
combining these gives the stated bound. \hfill$\square$

\begin{proposition}[Degeneration of divergence weighting]
\label{prop:divergence-collapse}
Let $\alpha_t^{\mathrm{div}} = \lvert \log \pi_\theta(y_t \mid y_{<t}, x)
- \log \pi_{\mathrm{ref}}(y_t \mid y_{<t}, x) \rvert$ and
$w^{\mathrm{div}} = (\alpha^{\mathrm{div}}+\varepsilon) /
\sum_t (\alpha_t^{\mathrm{div}}+\varepsilon)$.
Then
\[
0 \le \alpha_t^{\mathrm{div}} \le 2 L \beta \qquad \forall t,
\]
so if $\varepsilon$ is fixed and $\beta \to 0$, then
$w^{\mathrm{div}} \to 1/T$. If $\varepsilon=0$ or is negligible relative to
the raw scores, the same normalization is ill-conditioned as $\beta \to 0$:
the limiting weights are determined by tiny relative differences among
vanishing log-probability shifts rather than by a robust adapter-credit
signal.
\end{proposition}

The key observation is that the intrinsic scores are all determined by the
per-token \emph{log-probability differential}, which is precisely what LoRA's
rank-$r$ constraint drives small and approximately uniform across positions.
The degeneration is not a property of a particular weighting scheme; it is a
property of measuring per-token variation through the output distribution
when the policy has been constrained to a small neighborhood of the
reference.

\subsection{ARCA and Position-Discriminative Signal}\label{sec:theory-arca}

ARCA side-steps this degeneration by measuring a quantity that is position-varying
even when the adapter's logit-space impact is small.

\begin{proposition}[ARCA remains non-degenerate]
\label{prop:arca-nondeg}
Suppose the adapter-induced residual admits the decomposition $\Delta h_t =
\sum_\ell B_\ell A_\ell x_{\ell,t}^{\mathrm{pre}} + e_t$, where
$x_{\ell,t}^{\mathrm{pre}}$ is the pre-adapter activation at layer $\ell$ and
position $t$ and $\|e_t\|$ is a higher-order cross-layer term. Let
$v_{\ell,t} = B_\ell A_\ell x_{\ell,t}^{\mathrm{pre}}$ and define
$V_\ell = \mathrm{Var}_t(\|v_{\ell,t}\|)$ and $\bar{V}_\ell =
\mathbb{E}_t[\|v_{\ell,t}\|]$. Then
\[
\mathrm{Var}_t(\alpha_t^{\mathrm{ARCA}}) \ge \max_\ell V_\ell
- O\!\left(\sum_{\ell}\bar{V}_\ell \, \|e_t\|\right).
\]
Assume the residual scores dominate the normalization floor $\varepsilon$.
In particular, as long as the pre-adapter activations
$x_{\ell,t}^{\mathrm{pre}}$ have non-trivial position-to-position variance at
any layer, $\alpha_t^{\mathrm{ARCA}}$ has bounded-below variance across
positions, and its normalized weights $w_t^{\mathrm{ARCA}}$ do not converge
to the uniform distribution merely because the output-distribution shift is
small. With fixed $\varepsilon$, however, all scores become uniform in the
limit where the adapter residual itself vanishes below the floor.
\end{proposition}

The intuitive content of Proposition~\ref{prop:arca-nondeg} is that ARCA
inherits its position-discrimination from the \emph{base model's own
activation pattern} (which is already non-uniform because attention and
layer norm routing produce content-versus-filler distinctions) rather than
from a quantity that LoRA is specifically constrained to make small. Scaling
the adapter weights uniformly scales $\alpha^{\mathrm{ARCA}}$ uniformly, so
the normalized weights $w^{\mathrm{ARCA}}$ are scale-invariant in the
adapter magnitude.

\subsection{Bias--Variance Tradeoff}

Both advantage estimation and intrinsic token weighting introduce bias in
exchange for variance reduction. Advantage estimation does so by relying on
a learned value function, whose bias can be substantial when the underlying
state abstraction is weak. Intrinsic token weighting introduces bias by
reshaping the contribution of token-level gradients based on heuristic
salience measures. The empirical question is therefore not whether bias is
introduced, but whether the bias--variance tradeoff is favorable---and,
per Propositions~\ref{prop:surprisal-collapse}
and~\ref{prop:divergence-collapse}, whether the bias-inducing signal
survives the LoRA bottleneck at all.

\subsection{Summary}

This perspective reframes advantage estimation as one particular approach to
credit redistribution in policy gradient methods. Intrinsic token weighting
offers an alternative that leverages the structure of the language model's
predictive distribution, without assuming the existence of meaningful state
values. Under LoRA, however, the output-distribution-based signals that
most published weighting methods rely on can degenerate into uniform or
spuriously sparse token weights; the adapter-residual signal that ARCA uses
is, by Proposition~\ref{prop:arca-nondeg}, non-degenerate whenever the
adapter residual remains measurable above the normalization floor.
